%% it/sections/06_esiti.tex
%% ============================================================
%% 06_esiti.tex — Sistema punteggi ed esiti
%% ============================================================

\chapter{Gli Esiti / The Outcomes}
\markboth{Gli Esiti}{The Outcomes}

\lettrine[lines=3]{L}{a} sera del 30 marzo 1282 è l'inevitabile. La
domanda non è \textit{se} accadrà qualcosa — ma \textit{cosa}, e con
quale forma. Il sistema di punteggio misura quanto le fazioni hanno
preparato, organizzato, o ostacolato la notte che verrà.

\section{Il Sistema di Punteggio}

\subsection{Calcolo del Punteggio Totale}

Al termine della sessione, il GM coordinatore raccoglie i punteggi
delle quattro fazioni:

\begin{center}
\begin{tabular}{lc}
\toprule
\cinzel Fazione & \cinzel Punteggio (0--4) \\
\midrule
Congiurati   & \hspace{2cm} \\[4pt]
Baroni       & \hspace{2cm} \\[4pt]
Clero        & \hspace{2cm} \\[4pt]
Popolo       & \hspace{2cm} \\
\midrule
\textbf{Totale} & \textbf{\hspace{2cm} / 16} \\
\bottomrule
\end{tabular}
\end{center}

\nota{Per il GM Coordinatore}{
  Il punteggio viene calcolato a sessione conclusa, prima di rivelare
  l'esito ai giocatori. La rivelazione dell'esito è un momento narrativo:
  non limitatevi a leggere il testo. Adattatelo a ciò che è successo
  ai vostri tavoli.
}

\secsep

\section{Gli Esiti}

\subsection{0--4 punti — \textit{La Notte del Silenzio}}
\textcolor{oro}{\cinzel\large Fallimento della Cospirazione}

\filetto

La cospirazione è stata sgominata, o non era mai abbastanza organizzata
per funzionare. Jean de Saint-Rémy ha agito in anticipo: arresti
preventivi, coprifuoco, porto bloccato. La processione di Santo
Spirito si svolge sotto stretta sorveglianza militare.

Non accade nulla il 30 marzo.

Ma i semi sono stati piantati. Le ingiustizie restano. La rabbia
resta. I Vespri siciliani sono rimandati — non cancellati. La storia
del Mediterraneo sarà diversa. Quanto diversa, nessuno può saperlo.

\textbf{Conseguenze per fazione:}
\begin{itemize}[nosep]
  \item \textit{Congiurati} — la rete è parzialmente o totalmente
    smantellata. Giovanni da Procida dovrà ricominciare da capo.
  \item \textit{Baroni} — chi aveva preso posizione è compromesso.
    Chi non aveva preso posizione si congratula con sé stesso.
  \item \textit{Clero} — la Chiesa rimane al suo posto, né più
    né meno rispettata di prima.
  \item \textit{Popolo} — la repressione si abbatte sui quartieri.
    La gente comune paga il prezzo di ciò che altri hanno tentato.
\end{itemize}

\nota{Epilogo consigliato}{
  Leggete ai giocatori un breve testo ambientato un anno dopo: la
  rivolta è scoppiata altrove, in un modo diverso, con conseguenze
  imprevedibili. I loro personaggi sono vivi — o forse no. Ma la
  storia è andata avanti senza di loro.
}

\secsep

\subsection{5--8 punti — \textit{La Scintilla Spenta}}
\textcolor{oro}{\cinzel\large Rivolta Parziale e Repressa}

\filetto

La processione di Santo Spirito si svolge. C'è un incidente — una
perquisizione andata troppo oltre, un soldato che non sa quando
fermarsi. Si scatena una rissa. In alcuni quartieri la rissa diventa
violenza. In altri, le autorità riescono a contenere.

La rivolta c'è — ma è disorganizzata, localizzata, e viene soppressa
entro pochi giorni. Decine di morti tra la popolazione civile.
Carlo d'Angiò usa l'episodio per giustificare una repressione
più dura. Pietro d'Aragona non interviene: non era pronto, e la
finestra si è chiusa.

\textbf{Conseguenze per fazione:}
\begin{itemize}[nosep]
  \item \textit{Congiurati} — il piano era lì, ma non abbastanza
    robusto. Alcuni sono fuggiti. Alcuni sono stati catturati.
  \item \textit{Baroni} — chi aveva scommesso sulla rivolta perde
    tutto. Chi aveva tenuto il piede in due staffe sopravvive,
    a un prezzo.
  \item \textit{Clero} — viene aperta un'indagine ecclesiastica.
    L'arcivescovo è chiamato a rispondere a Roma.
  \item \textit{Popolo} — i quartieri piangono i morti. La memoria
    di questa notte alimenterà la prossima rivolta.
\end{itemize}

\secsep

\subsection{9--12 punti — \textit{Il Vespro}}
\textcolor{oro}{\cinzel\large Rivolta Completa}

\filetto

La processione di Santo Spirito diventa il detonatore. L'incidente
è reale, il momento è quello giusto, e la rete costruita dalle
fazioni tiene. In poche ore tutta Palermo è in rivolta.

I soldati francesi vengono cacciati o uccisi. Jean de Saint-Rémy
riesce a fuggire dal porto prima che venga bloccato. Entro
la mattina del 31 marzo, Palermo è libera per la prima volta
in sedici anni.

Questo è ciò che è accaduto storicamente — ma in questa versione,
i giocatori lo hanno costruito. Alcune cose sono andate diversamente.
Alcune persone sono vive che storicamente non lo erano, o viceversa.

\textbf{Conseguenze per fazione:}
\begin{itemize}[nosep]
  \item \textit{Congiurati} — la rete ha tenuto. Da Procida e
    Pietro d'Aragona hanno la loro opportunità. Ciò che ne faranno
    è un'altra storia.
  \item \textit{Baroni} — chi ha scommesso ha vinto — per ora.
    La guerra è appena cominciata.
  \item \textit{Clero} — le campane di Santo Spirito hanno suonato
    al momento giusto. La Chiesa siciliana ha una nuova autonomia —
    e nuove responsabilità.
  \item \textit{Popolo} — la vittoria appartiene a chi era in piazza.
    Ma le decisioni su cosa fare dopo non spetteranno a loro.
\end{itemize}

\nota{Varianti narrative}{
  Con 9--10 punti: la rivolta riesce, ma con più violenza e
  meno controllo di quanto previsto. Con 11--12 punti: la rivolta
  riesce in modo quasi preciso al piano. Adattate il racconto.
}

\secsep

\subsection{13--16 punti — \textit{L'Insurrezione}}
\textcolor{oro}{\cinzel\large Rivolta Totale e Coordinata}

\filetto

La notte del 30 marzo non è una rivolta — è un'insurrezione.
Le quattro fazioni hanno lavorato in modo abbastanza coordinato
da trasformare un incidente in una rivoluzione pianificata.
Non solo Palermo: i canali aperti verso Messina, Corleone,
Trapani hanno portato la rivolta a diffondersi in tutta la Sicilia
con una velocità che nessuno si aspettava — nemmeno i congiurati.

Pietro d'Aragona sbarca prima del previsto. La spedizione di Carlo
contro Bisanzio è cancellata permanentemente. Il Mediterraneo
è cambiato.

\textbf{Conseguenze per fazione:}
\begin{itemize}[nosep]
  \item \textit{Congiurati} — Giovanni da Procida è trionfante.
    Ma il trionfo ha un prezzo: chi decide cosa succede adesso
    è Pietro d'Aragona, non i siciliani.
  \item \textit{Baroni} — le casate che hanno scelto bene sono
    i nuovi pilastri del regno. Pesante responsabilità.
  \item \textit{Clero} — la Chiesa siciliana ha guadagnato
    un'autonomia senza precedenti — e un debito con Aragona.
  \item \textit{Popolo} — hanno vinto. I loro nomi non saranno
    nelle cronache. Ma loro lo sanno.
\end{itemize}

\nota{Nota di design — L'esito non è una ricompensa}{
  Nessuno degli esiti è intrinsecamente «buono» o «cattivo» in
  senso morale. Anche l'insurrezione totale ha un costo — la
  violenza della notte è reale, le conseguenze sono imprevedibili,
  e Pietro d'Aragona non è un liberatore disinteressato. Il gioco
  non premia chi ha fatto le scelte «giuste»: registra ciò che
  è successo, e lascia che i giocatori ne portino il peso.
}

\secsep

\section{Esiti Speciali per Fazione}

Indipendentemente dal punteggio totale, alcune condizioni speciali
modificano o aggiungono elementi alla narrazione finale:

\begin{itemize}
  \item \textbf{Congiurati a 0} — la rete è completamente smantellata,
    anche se la rivolta riesce: è scoppiata senza di loro, contro
    la loro volontà o in loro assenza.
  \item \textbf{Baroni a 4} — la nobiltà siciliana entra nella
    nuova era unita: raro e potente. Un personaggio baronale
    ottiene un titolo nella nuova Sicilia.
  \item \textbf{Clero a 4 e Sacrista vivo} — le campane di Santo
    Spirito suonano ancora mentre l'alba spunta. È il suono della
    vittoria. È il suono che ricorderanno.
  \item \textbf{Popolo a 4} — i quartieri popolari non subiscono
    nessuna ritorsione nella notte. I personaggi del Popolo e
    i loro cari sopravvivono tutti.
  \item \textbf{De Saint-Rémy Clock pieno} — anche in un esito
    di rivolta riuscita, la repressione ha fatto vittime. Leggete
    i nomi.
\end{itemize}

\filettodoppio

%% VERSIONE INGLESE


\filettodoppio